\scanpage{297}
\[
\begin{aligned}
&\deleted{=(\underline{\operatorname{Hom}}^*((\tau_NS))[-i],Q))_0}\\
&\deleted{=\underline{\operatorname{Hom}}^*(\tau_{N-i}S}\\
&\simeq(\underline{\operatorname{Hom}}^*(\tau_{N-i}(S),Q))_i\\
&\simeq\underline{\operatorname{Hom}}^i(\tau_{N-i}(S),Q)
\end{aligned}
\]

Enfin si dans $C$ les $\unsure{\underline{\lim}|\varinjlim}$ existent, alors $i_N$ admet un adjoint à gauche $q_N$.

En dualisant (appliquant ce qui précède à $C^o$, en utilisant l'isom

\diagram{\begin{tikzcd}
\mathcal M(C,S)^o\arrow[r,phantom,"\simeq"]\arrow[d]&\mathcal M(C^o,S)\arrow[d]\\
\mathcal M^N_{\deleted{\unsure{N}}}(C,S)^o\arrow[r,phantom,"\simeq"]&\mathcal M_{\leq N}(C^o,S)
\end{tikzcd}}

on voit que (si $\unsure{C,\ C^o}$ est une cat. de Groth. p ex) le foncteur
\[ i^N:\mathcal M^N\longrightarrow\mathcal M \]
\deleted{admet} \inserted{s'insère dans} une suite \deleted{\unsure{d'adj.}} \inserted{de 4} foncteurs adjoints
\[ \varphi^N\ \tau^N\ i^N\ q^N \]
qui se décrivent de façon duale.

\margin{298}
